\chapter{Document Authoring Guide}
\ChapterLabel{DocumentGuide}
\index{LaTeX}
\index{document conventions}
\index{Appendix!Document Authoring Guide}

\section{How to Use This Appendix}
Use this appendix to maintain the template's file conventions, build process,
and review rules. Update it when the document structure or toolchain changes.

\section{Template Provenance}
This template was developed by \textbf{Darian Muresan, Ph.D.}, in conjunction
with \textbf{OpenAI Codex} and \textbf{Anthropic Claude Code}. Keep this
attribution distinct from the assignment-specific instructor metadata.

\section{File Naming and Responsibilities}
Every chapter file shall use the \texttt{devops} prefix, a PascalCase chapter
name, and the \texttt{.tex} suffix. Each such file shall contain one complete
chapter, beginning with \texttt{\string\chapter\{...\}}, and all of its
sections and subsections. Do not split one chapter into section-level
\texttt{.tex} files. For example, all of Chapter~2 is contained in
\path{devopsExampleAssignment.tex}. Do not add requirements or
architecture-design chapters; this template records a DevOps learning activity
and its operational evidence.

\section{Synonyms and Controlled Terms}
\index{synonyms}
Use one preferred term for each important concept. If an assignment uses
equivalent terms, record them here before release. For example, state whether
``continuous delivery'' refers to \gls{cd} generally or specifically excludes
automatic production deployment. Add important terms to
\texttt{devopsGlossary.tex} and index them where used.

\section{Building the PDF}
The template requires a LaTeX distribution with \texttt{pdflatex},
\texttt{bibtex}, \texttt{makeglossaries}, and \texttt{makeindex}.
Code listings use the \texttt{minted} package, so LaTeX must also run with the
\texttt{-shell-escape} flag and Python Pygments must be installed
(\texttt{pip install Pygments}).
\subsection{Group Repository Build (Docker)}
\index{Docker!LaTeX build}
This group document is built with the TeX Live container described in
\ChapterRef{LatexDocker}, so no local LaTeX installation is required. From the
repository root, run:
\begin{verbatim}
docker build -t latex-builder latex-docker
docker run --rm -v "$PWD":/workdir -w /workdir/report \
  latex-builder devopsAssignment.tex
\end{verbatim}
\texttt{latexmk} reruns \texttt{pdflatex}, \texttt{bibtex},
\texttt{makeglossaries} (through \texttt{report/latexmkrc}), and
\texttt{makeindex} until all references settle. The PDF is written to
\texttt{report/devopsAssignment.pdf}. Review the log and visually inspect
affected PDF pages. Do not edit generated auxiliary files or the PDF directly.

\subsection{Local Toolchain Build}
A member who prefers a local toolchain can install the tools using the
platform instructions below and run the same \texttt{latexmk} command from the
\texttt{report} folder:
\begin{verbatim}
latexmk -pdf -shell-escape devopsAssignment.tex
\end{verbatim}

\section{Install LaTeX Tools on macOS}
\index{macOS}
\index{MacTeX}
\subsection{Recommended Installation}
Install the current full \href{https://www.tug.org/mactex/}{MacTeX}
distribution. It includes TeX Live, \texttt{pdflatex}, \texttt{bibtex},
\texttt{latexmk}, \texttt{makeglossaries}, \texttt{makeindex}, and useful
applications such as TeXShop and TeX Live Utility. The full distribution is the
least troublesome choice for this template because it includes the packages
used by the glossary, table, listing, color, and hyperlink features.

Either download the signed installer from the MacTeX web site and complete the
macOS installation dialog, or, if Homebrew is already installed, run:
\begin{verbatim}
brew install --cask mactex
\end{verbatim}
Close and reopen Terminal after installation. If the TeX commands are not
found, add the MacTeX binary directory to the current shell session:
\begin{verbatim}
export PATH="/Library/TeX/texbin:$PATH"
\end{verbatim}
Do not place this command in a shell profile until it has been verified on the
local computer.

\subsection{Verification and Compilation}
From the template directory, verify the tools:
\begin{verbatim}
command -v pdflatex bibtex latexmk makeglossaries makeindex
\end{verbatim}
Each command should print a path. Then create the PDF:
\begin{verbatim}
latexmk -pdf -shell-escape devopsAssignment.tex
\end{verbatim}
To remove generated intermediate files while keeping the PDF, run:
\begin{verbatim}
latexmk -c devopsAssignment.tex
\end{verbatim}

\section{Install LaTeX Tools on Windows}
\index{Windows}
\index{MiKTeX}
\subsection{Recommended Installation}
Install the current Basic MiKTeX distribution using the official
\href{https://miktex.org/howto/install-miktex}{MiKTeX Windows installation
guide}. During installation, choose a per-user installation when appropriate
and set missing-package installation to \emph{Ask me first} or \emph{Yes}.
The latter is convenient for a student machine with Internet access; the former
makes each package installation explicit.

Open MiKTeX Console after installation, apply available updates, and install
the \texttt{latexmk} package if the verification command below cannot find it.
On an organization-managed computer, obtain approval before installing software
or changing execution policy.

\subsection{Verification and Compilation}
Open a new PowerShell window in the template directory and verify:
\begin{verbatim}
Get-Command pdflatex,bibtex,latexmk,makeglossaries,makeindex
\end{verbatim}
Each command should display a command record. Build the PDF with:
\begin{verbatim}
latexmk -pdf -shell-escape devopsAssignment.tex
\end{verbatim}
If PowerShell blocks the script under a managed execution policy, ask the
system administrator for the approved procedure. Do not bypass institutional
security policy. To remove intermediate build files, run:
\begin{verbatim}
latexmk -c devopsAssignment.tex
\end{verbatim}

\section{Common Build Problems}
\begin{tabularx}{\linewidth}{L{3.0cm}Y}
\toprule \textbf{Symptom} & \textbf{Resolution}\\ \midrule
Command not found & Close and reopen the terminal. On macOS, verify the TeX
binary path; on Windows, verify MiKTeX completed installation and is on the user PATH.\\
Missing LaTeX package & Use MiKTeX Console or its approved on-the-fly package
installation prompt on Windows. On macOS, install the full MacTeX distribution
or use TeX Live Utility to add the missing package.\\
References or glossary incomplete & Run \texttt{latexmk} (or the Docker build) instead of
a single \texttt{pdflatex} command; it runs BibTeX, the glossary tool, and the
index tool in the required order.\\
minted error such as ``You must invoke LaTeX with the -shell-escape flag'' & Run
\texttt{pdflatex} with \texttt{-shell-escape}. If the error names Pygments or
\texttt{latexminted}, install Pygments with \texttt{pip install Pygments}.\\
Script is blocked & Follow the institution's approved PowerShell policy or run
the documented build commands with an approved administrator.\\ \bottomrule
\end{tabularx}

\section{Authoring and Review Rules}
Replace bracketed placeholders with assignment facts. Keep operational examples
safe, bounded, and free of credentials. Link activities to evidence, retain
history, and state unknowns or instructor decisions explicitly. Rebuild after
structural, cross-reference, table, glossary, or index changes.
